% Independent continuation by /root/frobenius.
\paragraph{C.4u (Proved: all supersingular threefolds in characteristic at least five).}
Every supersingular abelian threefold over an algebraically closed
field of characteristic \(p\ge5\) is an actual direct summand of a
product of Jacobians. In characteristics two and three we obtain,
respectively,
\[
       s_J(A)\mid2\quad(p=2),\qquad
       s_J(A)\mid3\quad(p=3);
\]
these two divisibilities do not by themselves prove a retraction.
More generally, for a supersingular abelian variety \(A\) of
dimension \(g\ge1\),
\begin{equation}\label{eq:c4-supersingular-p-bound}
       s_J(A)\mid p^{\,v_p((2g-1)!)}.
\end{equation}
In particular every such \(A\) is a direct summand if \(p>2g-1\).

\begin{proof}
We first justify the prime support of the required isogeny,
which matters integrally. In
\citet[Lemma 3.2 and its proof, author-version p.~7,
using Lemma 2.4, p.~5]{frC4YuSupersingular2017},
the supersingular variety is mapped to a superspecial variety
by replacing its covariant Dieudonn\'e lattice \(M\) by a
finite-index superspecial overmodule \(M'\) in the same
isocrystal. Our field is already perfect, so the purely
inseparable field extension in the statement is unnecessary.
The quotient \(M'/M\) has finite length over \(W(k)\) and is
killed by \(p^a\) for some \(a\ge0\). The associated isogeny
\[
                         \varphi:A\longrightarrow B
\]
therefore has group-scheme kernel killed by \(p^a\).
The superspecial variety \(B\) is a product of supersingular
elliptic curves over the algebraically closed field, hence
itself a product of Jacobians. This is the definition recalled
in the cited paper on p.~2; its Theorem 1.1 on p.~1 gives the
stronger uniqueness statement, which we do not need.

The quotient property supplies an actual homomorphism
\(h:B\to A\) with
\[
                            h\varphi=[p^a]_A.
\]
Consequently \(p^a\in I_J(A)\). Together with C.4i this shows
that \(s_J(A)\) is a power of \(p\) dividing \((2g-1)!\),
which proves \eqref{eq:c4-supersingular-p-bound}.
This argument works over arbitrary algebraically closed
fields of characteristic \(p\), and does not require that
\(A\) descend to a finite field.

For \(g=3\) the improved integral construction C.4p gives
\(6\in I_J(A)\). Thus
\[
                         s_J(A)\mid\gcd(6,p^a).
\]
This is one for \(p\ge5\), and divides two or three in the
two remaining characteristics. We can display the actual
retraction when \(p\ge5\). Choose
\(A\xrightarrow{u_6}J_6\xrightarrow{v_6}A\)
with \(v_6u_6=[6]_A\). Since \(p^{2a}\equiv1\pmod6\),
the integer \(c=(1-p^{2a})/6\) is defined. The maps
\[
 A\xrightarrow{\,(\varphi,u_6)\,}B\times J_6
       \xrightarrow{\,p^a h+c\,v_6\,}A
\]
have composite
\([p^{2a}+6c]_A=1_A\).
The splitting argument in C.4h therefore gives the required
product decomposition as abelian varieties.
\end{proof}

This settles the supersingular threefold case in every
characteristic at least five, including varieties without a
chosen principal polarization. All supersingular surfaces,
including characteristic two, were already covered by C.4s.
No assertion of principal polarizability of every supersingular
threefold is used here. The residual supersingular cases in
characteristics two and three, and the generic nonprincipal
surface cases discussed above, remain separate questions in
this investigation.
